\bookStart{Acre-Boot}[Æcer-bót]
\setBookCode{Acerbot}

\begin{flushright}%
\textbf{Dating:} ca. 1000

\textbf{Meter:} \Fornyrdislag%
\end{flushright}%

\section{Introduction}

The \textbf{Acre-Boot} (i.e.~‘Field-Cure’, abbrev. \Acerbot) is an unusually long text mixing Catholic, Germanic, and Celtic elements.

\Acerbot\ is preserved in Cotton MS Caligula A VII, foll.~176r--178r.

\sectionline

\section{Text}

\bpg\bpa Hér ys seo bót, hú þú meaht þíne æceras bétan gif hí nellaþ wel wexan oþþe þǽr hwilc un·ge·défe þing on ge·dón bið on dry oððe on lyb-lâce. Ge·nim þǫnne ǫn niht, ær hyt dagige, feower tyrf on feower healfa þæs landes, and ge·mearca hú hý ǽr stódon. Nim þonne ele and hunig and beorman, and ælces feos meolc þe on þæm lande sý, and ælces treow-cynnes dǽl þe on þǽm lande sý ge·wexen, bútan heardan béaman, and ælcre nam cúþre wyrte dǽl, bútan glappan ânon, and dó þonne hâlig-wæter ðǽr on, and drype þǫnne þriwa on þǫne staðol þára turfa, and cweþe ðonne ðás word:
„Crescite, wexe, et multiplicamini, and ge·mænig fealda, et replete, and ge·fylle, terre, þás eorðan.
In nomine patris et filii et spiritus sancti sit benedicti.“
And Pater Noster swá oft swá þæt óðer.  And bere siþþan ðá turf to circean, and mæsse-preost á·singe feower mæssan ofer þan turfon, and wende man þæt gréne to ðan weofode, and siþþan ge·bringe man þá turf þǽr hí ǽr wǽron ǽr sunnan setl-gange.
And hæbbe him gæ·worht of cwic-béame feower Cristes mælo and a·wríte on ælcon ende:
Matheus and Marcus, Lucas and Iohannes.
Lege þæt Cristes mæl on þone pyt neoþe-weardne, cweðe ðǫnne:
Crux Matheus, crux Marcus, crux Lucas, crux Sanctus Iohannes.\epa

\bpb TODO.\epb\epg


\bpg\bpa Nim ðǫnne þá turf and sete ðær ufon on and cweþe ðǫnne nigon síþon þás word:
Crescite, and swa oft Pater Noster, and wende þé þǫnne éast-weard, and on·lút nigon síðon éad-mód⸗líce, and cweð þǫnne þás word:\epa

\bpb TODO.\epb\epg


\bvg\bva%
„\alst{Éa}st-weard ic stande, \hld\ \alst{â}rena ic mé bidde, &
bidde ic þǫne \alst{m}ǽran domine, \hld\ bidde ðǫne \alst{m}iclan drihten, &
bidde ic ðǫne \alst{h}áligan \hld\ \alst{h}eofon-rices weard, &
\alst{eo}rðan ic bidde \hld\ and \alst{u}p-heofon &
and ðá \alst{s}óþan \hld\ \alst{s}ancta Marian &
and \alst{h}eofones meaht \hld\ and \alst{h}éah-reced, &
þæt ic móte þis \alst{g}ealdor \hld\ mid \alst{g}ife drihtnes &
\alst{t}óðum on·\alst{t}ýnan \hld\ þurh \alst{t}rumne ge·þanc, &
á·\alst{w}ęccan þas \alst{w}æstmas \hld\ ús to \alst{w}oruld-nytte, &
ge·\alst{f}ylle þás \alst{f}oldan \hld\ mid \alst{f}æste ge·léafan &
\alst{w}litigigan þas \alst{w}ancg-turf, \hld\ swá sé \alst{w}itega cwæð &
þæt sé hæfde \alst{á}re ǫn \alst{eo}rþ-ríce, \hld\ sé þe \alst{æ}lmyssan &
\alst{d}ǽlde \alst{d}óm-líce \hld\ \alst{d}rihtnes þances.“\eva

\bvb TODO.\evb\evg


\bpg\bpa Wende þe þǫnne III sun-ganges, a·strece þǫnne ǫn and-lang and a·rim þǽr letanias and cweð þonne: „Sanctus, sanctus, sanctus oþ ende.“\epa

\bpb TODO.\epb\epg


\bpg\bpa Sing þǫnne Benedicite a·þenedon earmon and Magnificat and Pater Noster III, and be·béod hit Criste and sancta Marian and þǽre hâlgan róde tó lofe and to weorþinga and þâm âre þe þæt land âge and eallon þâm þe him under·ðeodde synt.\epa

\bpb TODO.\epb\epg


\bpg\bpa Ðǫnne þæt eall síe ge·dón, þǫnne nime man un·cúþ sǽd æt ælmes-mannum and selle him twá swylc, swylce man æt him nime, and ge·gaderie ealle his sulh-ge·teogo to·gædere; borige þǫnne on þam béame stor and finol and ge·hâlgode sápan and ge·hâlgod sealt.
Nim þǫnne þæt sǽd, sete on þæs sules bodig, cweð þǫnne:\epa

\bpb TODO.\epb\epg


\bvg\bva%
„\alst{E}rce, \alst{E}rce, \alst{E}rce, \hld\ \alst{eo}rþan módor, &
ge·\alst{u}nne þé se \alst{a}l-walda, \hld\ \alst{é}ce drihten, &
æcera \alst{w}exendra \hld\ and \alst{w}rídendra, &
\alst{ea}cniendra \hld\ and \alst{e}lniendra; &
\alst{sc}eafta héhre \hld\ \alst{sc}íre wæstma, &
and þǽre \alst{b}râdan \hld\ \alst{b}ere-wæstma, &
and þǽre \alst{h}wítan \hld\ \alst{h}wæte-wæstma, &
and \alst{ea}lra \hld\ \alst{eo}rþan wæstma.\eva

\bvb TODO.\evb\evg


\bvg\bva%
Ge·\alst{u}nne him \hld\ \alst{é}ce drihten &
and his \alst{h}âlige, \hld\ þe ǫn \alst{h}eofonum synt, &
þæt hys yrþ sí ge·\alst{f}riþod \hld\ wið ealra \alst{f}eonda ge·hwæne, &
and heo sí ge·\alst{b}orgen \hld\ wið ealra \alst{b}ealwa ge·hwylc, &
þára \alst{l}yb-lâca \hld\ geond \alst{l}and sáwen.\eva

\bvb TODO.\evb\evg


\bvg\bva%
Nú ic bidde ðǫne \alst{w}aldend, \hld\ sé ðe ðás \alst{w}or-uld ge·sceóp, &
þæt ne sý nán tó þæs \alst{c}widol wíf \hld\ ne tó þæs \alst{c}ræftig man &
þæt á·\alst{w}endan ne mæge \hld\ \alst{w}ord þus ge·cwedene.“\eva

\bvb TODO.\evb\evg



\bpg\bpa Þonne man þá sulh forð drife þá forman furh on·scéote, cweð þǫnne:\epa

\bpb TODO.\epb\epg


\bvg\bva%
„Hál wes þú, \alst{F}olde, \hld\ \alst{f}ir\emph{h}a módor! &
Beo þú \alst{g}rówende \hld\ on \alst{G}odes fæþme, &
\alst{f}ódre ge·\alst{f}ylled \hld\ \alst{f}ir\emph{h}um tó nytte.“\eva

\bvb “Hale be thou, Fold, mother of men! \\
Be thou growing in God’s bosom, \\
filled with fodder as a boon to men.”\evb\evg


\bpg\bpa Nim þonne ælces cynnes melo and a·bacæ man inn-werdre handa brádnæ hláf and ge·cned hine mid meolce and mid hálig-wætere and lecge under þa forman furh. Cweþe þonne:\epa

\bpb TODO.\epb\epg


\bvg\bva%
„\alst{F}ul æcer \alst{f}ódres \hld\ \alst{f}ir\emph{h}a cinne, &
\alst{b}eorht-\alst{b}lowende, \hld\ þú ge·\alst{b}létsod weorþ &
þæs \alst{h}âligan nǫman \hld\ þe ðás \alst{h}eofon ge·sceop &
and þás \alst{eo}rþan \hld\ þe wé \alst{ǫ}n lifiaþ; &
sé \alst{g}od, sé þás \alst{g}rundas ge·worhte, \hld\ ge·unne ús \alst{g}rowende gife, &
þæt ús \alst{c}orna ge·hwylc \hld\ \alst{c}ume tó nytte.“\eva

\bvb “Acre, full of fodder for the race of men, \\
brightly blooming; thou wast blessed \\
by the holy name which created the heaven \\
and this earth upon which we live. \\
The God who made the lands, grant us growing grace \\
that each kind of grain may be for us a boon.”\evb\evg


\bpg\bpa Cweð þonne: „III Crescite in nomine patris, sit benedicti.“  Amen and Pater Noster þríwa.\epa

\bpb TODO.\epb\epg
